\bookchapter{الفصل 33: نظرة مستقبلية — المحطة القادمة للهاتف}

\chapterintro{لن يختفي الهاتف بين عشية وضحاها، لكن “سيادته” قد تُقسم لأجهزة الجيل الجديد. هذا الفصل نتحدث عن ثلاثة: كيف تبدو الأجهزة القادمة، كيف يتحول الهاتف لـ“مخفي” كأساس للهوية، والقضايا الأخلاقية التي يجب التفكير فيها مسبقاً.}

\booksection{واحد: AR/VR والحوسبة المكانية: “سطح المكتب الثاني” خارج الهاتف}

منافسة الجيل القادم من الأجهزة بدأت، والأكثر احتمالاً لتقسيم وقت الهاتف هي أجهزة AR/VR والحوسبة المكانية:

\begin{itemize}
  \item \textbf{خلف Vision Pro}: إطلاق Apple Vision Pro يدل على هبوط مفهوم "الحوسبة المكانية"، لكن الجيل الأول سعره (29999 يوان فأكثر) ووزنه لا يزالان عائقاً. يُتوقع 2026-2027 إطلاق نسخة أخف وأرخص، لينتقل من "لعبة المطور" لـ"الجمهور المبكر". حينها قد يصبح iPhone رفيق الحوسبة ومركز البيانات لنظارة Vision.
  \item \textbf{من "مركز الأجهزة" إلى "مركز الهوية"}: افتح "سطح مكتب ثانياً" أمام عينيك، المعلومات الرقمية تصبح كالهواء في كل مكان. نظارات Meta وRay-Ban الذكية باعت ملايين الأزواج، والشركات المحلية Rokid وXREAL وINMO تطلق نظارات AR استهلاكية. نظارات AR 2025-2026 تترجم فورياً وتعرض الملاحة وتتحدث مع مساعد AI، لكنها لا تزال محدودة بعرض دليل الموجة والبطارية، وتبعد 3-5 سنوات عن "استبدال الهاتف".
  \item \textbf{الخواتم الذكية والقابلة للارتداء}: Samsung Galaxy Ring وOura Ring وغيرها من الخواتم الذكية بدأت تنتشر، كمستشعر "ملاصق" للهاتف يكمل مراقبة الصحة. دور الهاتف يتحول من "طرفي شامل" لـ"مركز + عدة ملحقات خفيفة".
\end{itemize}

\booksection{اثنان: واجهة الدماغ والحاسوب: الفكر كوسيلة تفاعل}

الفكر نفسه سيصبح وسيلة للتفاعل، لكن هذا الطريق أطول مما تتخيل:

\begin{itemize}
  \item \textbf{أحدث تطورات Neuralink}: Neuralink لإيلون ماسك أكملت 2024 أول زرع بشري، أول مشارك Noland Arbaugh يتحكم في ماوس الحاسوب بالتفكير ويلعب الشطرنج. حالياً BCI (واجهة الدماغ والحاسوب) الزرعي يوجه أساساً لمساعدة مرضى الشلل، الهدف مساعدتهم على "تشغيل" العالم الرقمي من جديد.
  \item \textbf{BCI غير الجراحي}: عدة شركات تطور أطواق رأس غير جراحية (مثل Kernel وOpenBCI)، تقرأ إشارات الدماغ دون جراحة، لكن الدقة أقل بكثير من الزرعي. يبعد كثيراً عن "الكتابة بالتفكير".
  \item \textbf{التوقعات الاستهلاكية}: واجهة الدماغ والحاسوب تبعد 10-20 سنة عن التفاعل الاستهلاكي مع الهاتف، لكن قصيرة المدى قد تُستخدم أولاً في المساعدة للوصول ومراقبة التركيز وغيرها. قضايا الأخلاقيات والخصوصية ("موجات دماغي لمن تعتبر خصوصية") ستصبح نقطة جدال قبل نضج التقنية.
\end{itemize}

\booksection{ثلاثة: “اختفاء” الهاتف وأساس الهوية الجديد}

اتجاه تطور آخر للهاتف قد يكون “الاختفاء” — التحمّل كـ“محرك خفي” في كل مكان، ليصبح الأساس الوحيد لهويتك الرقمية:

\begin{itemize}
  \item \textbf{رابعًا: حافظ على الاتصال، وحافظ على الاستقلالية — فلسفتك للحياة الرقمية}: مصادقة "أنت أنت" فورياً، تسجيل الدخول لأي خدمة دون إدخال رمز تحقق يدوياً.
  \item \textbf{اجعل الأتمتة "مساعدك الثاني" الموثوق.}: عند تنشيط جهاز جديد، إعداد مسبق لخطة ميراث الأصول الرقمية، لجعل "ما بعد الوفاة" افتراضياً من المصنع.
  \item \textbf{نفذ بانتظام "تنقيحًا رقميًا" و"صومًا اجتماعيًا".}: خزنة بياناتك الشخصية السحابية هي المركز، والعتاد مجرد مدخل وصول. تغيير الجهاز، الهوية والبيانات تتبعك دائماً.
\end{itemize}

\booksection{أربعة: تأملات أخلاقية: هل نحن سادة الهاتف؟}

اكتب “شروط خدمة” بسيطة لك وللهاتفك لتحديد علاقتكما. مثل:

\begin{itemize}
  \item \textbf{أخلاقيات AI}: عندما يفهم AI الجانبي تفضيلاتك أفضل منك، هل "اختيارك" إرادة حرة أم نتيجة توجيه؟ عندما يطلب AI Agent ويرد الرسائل بالنيابة عنك، أين المسؤولية والحدود؟
  \item \textbf{سيادة البيانات}: بياناتك الصحية وموجات الدماغ والبيانات البيومترية، لك أم لشركة الأجهزة أم لمزود السحابة؟ قانون حماية المعلومات الشخصية وGDPR الأوروبي مجرد بداية، و"حق نقل البيانات" و"الحق في النسيان" ما زال طريقهما طويلاً للتطبيق.
  \item \textbf{الفجوة الرقمية}: هواتف AI والشاشات القابلة للطي والاتصال الفضائي توسع الفجوة بين "يستطيع الشراء" و"لا يستطيع". الوصول العادل للبنية التحتية الرقمية سيصبح قضية اجتماعية جديدة — مثل "هل تعرف استخدام الإنترنت" قبل 20 سنة.
  \item البند الأول: وقت انتباهي ملكي. يجب أن تدخل تلقائيًا حالة "عدم الإزعاج" عندما أعمل/أرافق عائلتي.
  \item البند الثاني: بياناتي الشخصية هي ملكي. واجبك الأساسي هو التشفير الآمن والمزامنة المريحة.
\end{itemize}

\booksection{خمسة: ابقَ متصلاً وابقَ مستقلاً — فلسفة حياتك الرقمية}

مهما تطورت التقنية، فلسفة الحياة الرقمية التالية صالحة بعد 10 سنوات:

\begin{enumerate}
  \item خامساً: الملحق
  \item حسنًا، بعد "رحلة طويلة" عبر أكثر من عشرين فصلاً، تهانينا! لقد تطورت بنجاح من "مبتدئ في الهواتف" وأتقنت المهارات الأساسية من المبتدئ إلى المتقدم. هذا الملحق يشبه دليل مرجعي سريع محمول وصندوق كنوز موارد لك.
  \item الملحق A: عملية فتح صندوق الهاتف الجديد "خطوة بخطوة"
  \item 1. فحص الصندوق "الثلاثي":
  \item فحص المظهر
\end{enumerate}

\begin{tcbpractice}{:افحص جسم الهاتف والشاشة بحثًا عن خدوش أو اصطدامات.}
اكتب “شروط خدمة” بسيطة لك ولهاتفك (ونظارة AR ومساعد AI المستقبلية)، لتعريف علاقتكما. مثلاً:

\begin{itemize}
  \item :تحقق من مطابقة IMEI/الرقم التسلسلي على الصندوق مع الفاتورة.
  \item اختبار الوظائف
  \item :أجرِ مكالمة لاختبار سماعة وميكروفون، والتقط صورة لفحص الكاميرا.
\end{itemize}

\end{tcbpractice}

\begin{tcbsummary}{خلاصة الفصل}
\begin{itemize}
  \item Vision Pro ونظارات AR مدخل حوسبة جديد خارج الهاتف، لكن صعب استبداله قصيرة المدى
  \item واجهة الدماغ والحاسوب تبعد 10-20 سنة عن الاستهلاك، القضايا الأخلاقية قبل نضج التقنية
  \item الهاتف يتطور نحو "الاختفاء"، من مركز الأجهزة لمركز الهوية
  \item أخلاقيات AI وسيادة البيانات والفجوة الرقمية القضايا الأساسية للعشر سنوات القادمة
  \item ابقَ متصلاً وابقَ مستقلاً، فلسفة الحياة الرقمية الثابتة
\end{itemize}

\end{tcbsummary}
